\itemhead{3.2 例题 · 例题5（教材p.84）}{XJ-B1-C03-EX-P084-Y0592}
\label{XJ-B1-C03-EX-P084-Y0592}
{\footnotesize 湘教版教材p.84 / PDF第91页 / 3.2 例题 / 例题5}
\field{原题}{设 $g(x)$ 是定义于 $[-5,5]$ 上的函数，且 $f(x)=g(x)+g(-x)$，讨论 $f(x)$ 的奇偶性；如果在 $[0,5]$ 上 $f(x)=1-2x$，试求 $f(x)$ 在 $[-5,0]$ 上的解析式。}
\field{参考答案}{$f(-x)=g(-x)+g(x)=f(x)$，故为偶函数；$x\in[-5,0]$ 时 $f(x)=f(-x)=1+2x$。}
{\small \field{关联知识}{F08 奇偶性、对称与函数性质应用（教材帮p.113、114、115 / PDF122、123、124）；F03 函数求值、解析式与函数建模（教材帮p.80、87、88、89、90、131 / PDF89、96、97、98、99、140）}
\field{方法与易错点}{检查定义域，再代入 $-x$；求另一侧解析式时同时写出 $-x$ 所在区间。 易错：只看公式不看定义域；用一对点证明奇偶；负半轴代换遗漏最外层负号。}
\field{答案依据}{例题原解位置：湘教版教材p.84及题库原图；独立编写可编辑参考解答。答案册此部分主要收练习、习题，不伪造例题答案册位置。}}
